% Generated draft: requires mathematical and source review.
\documentclass[11pt]{article}
\usepackage[margin=25mm]{geometry}
\usepackage{fontspec}
\setmainfont{Latin Modern Roman}
\usepackage{amsmath,amssymb,mathtools,unicode-math}
\setmathfont{Latin Modern Math}
\usepackage{enumitem,needspace,xurl,hyperref}
\setlength{\parindent}{0pt}
\setlength{\parskip}{5pt}
\setlength{\emergencystretch}{3em}
\newcommand{\N}{\mathbb{N}}
\newcommand{\Z}{\mathbb{Z}}
\newcommand{\Q}{\mathbb{Q}}
\newcommand{\R}{\mathbb{R}}
\newcommand{\C}{\mathbb{C}}
\newcommand{\F}{\mathbb{F}}
\newcommand{\A}{\mathbb{A}}
\newcommand{\PP}{\mathbb{P}}
\newcommand{\E}{\mathbb{E}}
\newcommand{\eps}{\varepsilon}
\newcommand{\dd}{\,\mathrm{d}}
\begin{document}
\section*{UnsolvedMath ID: 800982}
\textbf{Source problem code:} catalogue entry
\subsection{Definitions and mathematical statement}
Let \((M,g)\) be a closed Riemannian surface, meaning a compact Riemannian surface without boundary. No orientability or connectedness hypothesis is stated in the supplied record.

The supplied excerpt does not define the admissible class of stationary nets. In the usual unweighted geodesic-network description, edges are geodesic arcs and stationarity at each vertex \(v\) requires the balance condition
\[
\sum_{e\text{ incident at }v}T_{e,v}=0\in T_vM,
\]
where \(T_{e,v}\) is the unit tangent pointing away from \(v\) along the corresponding incident edge germ. The valency \(\deg_N(v)\) counts incident edge germs. This describes the usual local stationarity condition; the supplied excerpt does not determine all global admissibility conventions, including finiteness, weights, or permitted intersections.

Let \(\mathcal S(M,g)\) denote the class of stationary nets intended by the source, with \(|N|\subseteq M\) the image of a net \(N\). Following the source, call \(N\) non-trivial if
\[
\exists v\in V(N)\qquad \deg_N(v)\geq 3.
\]
The source also describes this as saying that \(N\) is not a disjoint union of simple closed geodesics; the equivalence depends on the precise admissibility and vertex conventions for stationary nets.

The question attributed to Gromov is whether, for every closed Riemannian surface \((M,g)\),
\[
\forall U\subseteq M\text{ open},\quad U\neq\varnothing
\quad\Longrightarrow\quad
\exists N\in\mathcal S(M,g):\quad
\bigl(\exists v\in V(N),\ \deg_N(v)\geq3\bigr)
\ \text{and}\ |N|\cap U\neq\varnothing.
\]
Equivalently, is
\[
\overline{\bigcup_{\substack{N\in\mathcal S(M,g)\\N\text{ non-trivial}}}|N|}=M?
\]
Here density refers to the union of the images of all non-trivial stationary nets. The source's final reformulation omits the word ``non-trivial''; the displayed formulation retains the restriction from its main question. A fully specified formulation requires confirmation of the intended stationary-net conventions.
\subsection{Short English statement}
On every closed Riemannian surface, must every nonempty open set meet a stationary net having a vertex of valency at least three? This asks whether the images of all such nets together are dense. The supplied excerpt leaves the precise admissibility conventions for stationary nets unspecified.
\subsection{Sources}
\begin{itemize}
\item Local UnsolvedMath registry, permanent ID \(800982\), source problem code \texttt{catalogue entry}; supplied complete source record and provenance.
\item Burns and Matveev, \emph{Open Problems and Questions about Geodesics} (2019), Question 63, attributed in the supplied record to Gromov. Supplied source URL: \url{https://arxiv.org/abs/1308.5417}.
\item The supplied provenance records access on 2026-10-16 and extraction by \texttt{arxiv\_latex}. These are supplied provenance details; no independent source visit or verification is asserted here.
\item The registry labels the problem ``open''. This is a supplied status label, not an independently verified mathematical or literature assessment.
\end{itemize}
\end{document}
